\documentclass[10pt,a4paper]{amsart}
\usepackage[margin=19mm]{geometry}
\usepackage{amsmath,amssymb,mathtools}
\usepackage[scr=rsfs,cal=euler]{mathalfa}
\usepackage{xfrac}
\usepackage[unicode,hidelinks,bookmarks=false]{hyperref}
\newcommand{\ie}{i.e.\ }
\newcommand{\cf}{cf.\ }
\newcommand{\resp}{respectively\ }
\newcommand{\Subsection}{\subsection}
\providecommand{\nobreakdash}{\nobreak\hbox{-}}
\setlength{\emergencystretch}{2em}
\multlinegap0pt
\usepackage{mathtools}
\usepackage{amssymb}
\usepackage{xcolor}
\newtheorem{theorem}{Theorem}
\newtheorem{proposition}{Proposition}
\newcommand{\C}{\mathcal C}
\newcommand{\D}{\mathcal D}
\newcommand{\Hh}{\mathcal H}
\newcommand{\Irr}{\operatorname{Irr}}
\newcommand{\M}{\mathcal M}
\newcommand{\N}{\mathcal N}
\DeclareMathOperator{\id}{id}
\title{Derived Reconstruction of Semisimple Tensor and Module Categories}
\author{Sheng Tan}
\date{Source: arXiv:2609.06753v1 (CC BY 4.0)}
\begin{document}
\maketitle

\noindent\emph{Source and licence.} Derived Reconstruction of Semisimple Tensor and Module Categories. Version arXiv:2609.06753v1 (CC BY 4.0), distributed by arXiv under the Creative Commons Attribution 4.0 International licence, http://creativecommons.org/licenses/by/4.0/, as recorded in the arXiv licence metadata for this exact version and shown on the abstract page https://arxiv.org/abs/2609.06753v1. Full licence notice: \texttt{licenses/arxiv-2609.06753-CC-BY-4.0.txt}.\par\smallskip

\noindent\textit{Editorial scope.} Counted unit: the article's theorem thm:simultaneous-reconstruction (source0.tex lines 643--645). The article's complete original proof of it is delivered with it (source0.tex lines 647--697, one \texttt{proof} environment of 51 lines). No mathematics is written, repaired or re-derived: the statement and the proof are the article's own lines, in the article's own notation, and no retained line is substituted or re-flowed.  Retained uncounted as the source-local context the proof needs: lines 594--600; lines 627--633.  Nothing was omitted from any retained span, so the package carries no omission record.  No retained line contains a citation, so the wrapper carries no bibliography and no bibliography entry.  No retained line contains a figure environment, an image-inclusion command or drawing code, so no image is delivered and the package declares no asset dependency.  Editorial formatting only, in addition: an AMS \texttt{amsart} preamble that restates the article's own declarations and macros used by the retained lines, namely the environments \texttt{theorem}, \texttt{proposition} on separate counters, together with the standard packages those lines need.  The retained environments print in this document as Theorem 1, Proposition 1 in order of appearance, whereas the article prints the same items with the numbering of its own full document.  The complete native source of the article as distributed in the arXiv e-print for this exact version is bundled unchanged as \texttt{sources/209566/source0.tex}.
\begin{theorem}\label{thm:module-regrading}
The functor
\[
R_p:\M\rightarrow\Hh_p,\quad R_p(M)=\bigoplus_r M_r[r],
\]
is an exact equivalence. It admits a coherent $R_\chi$-module-functor structure and is an $R_\chi$-module equivalence.
\end{theorem}

\begin{proposition}\label{prop:transported-pair}
Let $F: D^b(\C)\to D^b(\D)$ be a monoidal triangulated equivalence, and let $G: D^b(\M)\to D^b(\N)$ be a triangulated equivalence with a coherent $F$-module-functor structure. Transport the standard $t$-structures along $F$ and $G$. The resulting $t$-structure on $D^b(\D)$ is associated with a unique character $\chi: U_{\D}\to\mathbb Z$. Let $p:\Irr(\N)\to\mathbb Z$ be the shift function of the resulting $t$-structure on $D^b(\N)$. Then
\[
p(P)=p(M)+\chi(\deg X)
\]
whenever $X\in\Irr(\D)$, $M,P\in\Irr(\N)$, and $P$ is a direct summand of $X\odot M$.
\end{proposition}

\begin{theorem}\label{thm:simultaneous-reconstruction}
Let $\C$ and $\D$ be essentially small locally finite semisimple tensor categories, and let $\M$ and $\N$ be nonzero essentially small locally finite semisimple left module categories over $\C$ and $\D$, respectively. Suppose that $F: D^b(\C)\to D^b(\D)$ is a monoidal triangulated equivalence and that $G: D^b(\M)\to D^b(\N)$ is a triangulated equivalence with a coherent $F$-module-functor structure. Then there are an exact $k$-linear tensor equivalence $f:\C\to\D$ and an exact $k$-linear $f$-module equivalence $g:\M\to\N$.
\end{theorem}

\begin{proof}
Let $\Hh_\chi$ and $\Hh_p$ be the hearts of the transported $t$-structures from Proposition~\ref{prop:transported-pair}. Restriction gives exact equivalences
\begin{equation*}
\overline F=F|_{\C}:\C\rightarrow\Hh_\chi,\quad \overline G=G|_{\M}:\M\rightarrow\Hh_p.
\end{equation*}
The tensorator of $F$ makes $\overline F$ strong monoidal, and the module constraint of $G$ makes $\overline G$ an $\overline F$-module equivalence. Let
\begin{equation*}
R=R_\chi:\D\rightarrow\Hh_\chi,\quad P=R_p:\N\rightarrow\Hh_p
\end{equation*}
be the equivalences constructed above. Write $\mu_{X,Y}: R(X)\otimes R(Y)\to R(X\otimes Y)$ for the tensorator of $R$, and let $\rho_{X,M}: R(X)\odot P(M)\to P(X\odot M)$ be the isomorphism from Theorem~\ref{thm:module-regrading}.

Choose quasi-inverses $S:\Hh_\chi\to\D$ and $Q:\Hh_p\to\N$, with counits
\begin{equation*}
\varepsilon: RS\xrightarrow{\sim}\id_{\Hh_\chi},\quad
\delta: PQ\xrightarrow{\sim}\id_{\Hh_p}.
\end{equation*}
We first give $S$ a monoidal structure. Define the tensorator $S(A)\otimes S(B)\to S(A\otimes B)$ as the unique map whose image under $R$ satisfies
\begin{equation}\label{eq:S-tensorator}
\varepsilon_{A\otimes B}\circ R\bigl(S(A)\otimes S(B)\to S(A\otimes B)\bigr)\circ\mu_{S(A),S(B)}
=\varepsilon_A\otimes\varepsilon_B.
\end{equation}
Define the unit map in the same way. Since $R$ is fully faithful, these maps are uniquely determined. Naturality follows from that of $\mu$ and $\varepsilon$. The coherence of $\mu$ then gives the associativity and unit conditions. Thus $S$ is strong monoidal and $\varepsilon$ is monoidal.

We next give $Q$ an $S$-module structure. For $A\in\Hh_\chi$ and $B\in\Hh_p$, define
\begin{equation*}
\kappa_{A,B}: S(A)\odot Q(B)\rightarrow Q(A\odot B)
\end{equation*}
as the unique map satisfying
\begin{equation}\label{eq:kappa}
\delta_{A\odot B}\circ P(\kappa_{A,B})\circ\rho_{S(A),Q(B)}
=\varepsilon_A\odot\delta_B.
\end{equation}
Such a map exists and is an isomorphism because $P$ is fully faithful. Its naturality follows from that of $\rho$, $\varepsilon$, and $\delta$. Applying $P$ to the module associativity diagram and using \eqref{eq:kappa} reduces it to the associativity diagram for $\rho$ together with the monoidality of $\varepsilon$. The unit condition follows in the same way. Hence $\kappa^{-1}$ gives a coherent module-functor structure on $Q$.

Finally, invert the restricted module constraint of $\overline G$ and write it as
\begin{equation*}
\gamma_{X,M}:
\overline F(X)\odot\overline G(M)
\xrightarrow{\sim}
\overline G(X\odot M).
\end{equation*}
Set $f=S\overline F$ and $g=Q\overline G$. Then $f$ is a tensor equivalence. Define
\begin{equation*}
\beta_{X,M}
Q(\gamma_{X,M})\circ
\kappa_{\overline F(X),\overline G(M)}
:
f(X)\odot g(M)\rightarrow g(X\odot M).
\end{equation*}
The associativity condition for $\beta$ follows from those for $\kappa$ and $\gamma$, together with the monoidal structure on $f$. The unit condition follows in the same way. Thus $\beta^{-1}$ gives a coherent $f$-module-functor structure on $g$. Therefore $g$ is an $f$-module equivalence. Since $f$ and $g$ are equivalences between semisimple abelian categories, they are exact.
\end{proof}

\end{document}
